\subsection*{Week of 2026-11-23: sums, integrals and limits}

Second rail: inequalities. The first timed half paper is on Saturday 2026-11-28.

\begin{center}\small
\begin{tabular}{@{}c >{\raggedright\arraybackslash}p{3.1cm} >{\raggedright\arraybackslash}p{4.8cm} >{\raggedright\arraybackslash}p{4.8cm}@{}}
\toprule
Rung & Problem & Object and ask & Move \\
\midrule
1 & \seminar{Limits \#15} and \seminar{Limits \#12} & A Riemann sum and a product limit from older Putnams. & Recognise an integral in a sum, and logarithms to turn a product into a sum. \\
2 & \putnam{2016 B1} & A sequence defined through logarithms; sum a series from it. & Telescoping. \\
3 & \putnam{2022 A1} & A line and a logarithmic curve meeting exactly once; the parameters. & Convexity and tangency. \\
4 & \putnam{2019 B2} & A trigonometric sum; its asymptotic limit. & A sum that is a Riemann sum in disguise, with a correction term. \\
5 & \putnam{2021 A2} & A function defined by a limit; an asymptotic ratio. & Take logarithms, expand, keep the first surviving term. \\
6 & \putnam{2025 A2} & A sine bounded by a quadratic; the best constants. & Tagged inequality; the limit at zero and the equality case carry it. Cite the rung where you first expanded to the first surviving term. \\
\midrule
opt & \putnam{2020 A2} & A binomial sum; closed form. & \\
opt & \putnam{2024 B3} & Gaps between fixed points of a tangent-type function; a two-sided bound. & \\
opt & \putnam{2023 A1} & A product of cosines; its second derivative exceeds a threshold. & \\
\bottomrule
\end{tabular}
\end{center}

Seminar: \seminar{PS 2 \#1, \#5, \#8, \#15, \#16}; \seminar{PS 5 \#3, \#4, \#13}; \seminar{Limits \#1} to \seminar{\#4}, \seminar{\#11}, \seminar{\#13}, \seminar{\#14}, \seminar{\#16}.

\textbf{Mock 1, Saturday 2026-11-28.} The 2015 A paper as two sittings of 90 minutes with fifteen minutes between them, the shape of sessions A and B: A1, A3 and A5 in the first sitting, A2, A4 and A6 in the second. The split gives each sitting an approachable, a middle and a hard problem; that is a guess at how the new sessions are built, not something the policy says. Start at 11:00 and 12:45 if the day allows, the Eastern start times of sessions A and B. Room rules as in \secref{sec:examday}{exam day}: no phone or watch, one side of the paper, one problem per sheet, pens down at 90 minutes. Three call lines go on the scratch paper in the first eight minutes of each sitting. In the break, before looking at the scratch sheet again, write the three lines from memory, then compare; on the day the scratch paper is handed in, so memory is all that leaves the room. Type the six lines in that evening with the date. Settle every one, including those not attempted; \emph{moot} is the honest verb for those, and seeing how many there are is part of the point.
